\documentclass[10pt,a4paper]{amsart}
\usepackage[margin=19mm]{geometry}
\usepackage{amsmath,amssymb,mathtools}
\usepackage[scr=rsfs,cal=euler]{mathalfa}
\usepackage{xfrac}
\usepackage[unicode,hidelinks,bookmarks=false]{hyperref}
\newcommand{\ie}{i.e.\ }
\newcommand{\cf}{cf.\ }
\newcommand{\resp}{respectively\ }
\newcommand{\Subsection}{\subsection}
\providecommand{\nobreakdash}{\nobreak\hbox{-}}
\setlength{\emergencystretch}{2em}
\multlinegap0pt
\usepackage{amssymb}
\usepackage{xfrac}
\usepackage{tikz}
\usepackage[shortlabels]{enumitem}
\newtheorem{lemma}{Lemma}
\def\ZZ{{\mathbb Z}}
\newcommand{\s}{\sigma}
\title{Emergence of a convex strange term via homogenization \\ of non-local energies at the critical exponent}
\author{Giuseppe Cosma Brusca, Giuliana Fusco}
\date{Source: arXiv:2609.11371v1 (CC BY 4.0)}
\begin{document}
\maketitle

\noindent\emph{Source and licence.} Emergence of a convex strange term via homogenization \\ of non-local energies at the critical exponent. Version arXiv:2609.11371v1 (CC BY 4.0), distributed by arXiv under the Creative Commons Attribution 4.0 International licence, http://creativecommons.org/licenses/by/4.0/, as recorded in the arXiv licence metadata for this exact version and shown on the abstract page https://arxiv.org/abs/2609.11371v1. Full licence notice: \texttt{licenses/arxiv-2609.11371-CC-BY-4.0.txt}.\par\smallskip

\noindent\textit{Editorial scope.} Counted unit: the article's lemma cammini (source0.tex lines 1926--1936). The article's complete original proof of it is delivered with it (source0.tex lines 1938--2088, one \texttt{proof} environment of 151 lines). No mathematics is written, repaired or re-derived: the statement and the proof are the article's own lines, in the article's own notation, and no retained line is substituted or re-flowed.  No further source-local context is needed: every cross-reference of every retained line resolves inside the two delivered spans.  Nothing was omitted from any retained span, so the package carries no omission record.  No retained line contains a citation, so the wrapper carries no bibliography and no bibliography entry.  No retained line contains a figure environment, an image-inclusion command or drawing code, so no image is delivered and the package declares no asset dependency.  Editorial formatting only, in addition: an AMS \texttt{amsart} preamble that restates the article's own declarations and macros used by the retained lines, namely the environments \texttt{lemma} on separate counters, together with the standard packages those lines need.  The retained environments print in this document as Lemma 1 in order of appearance, whereas the article prints the same items with the numbering of its own full document.  The complete native source of the article as distributed in the arXiv e-print for this exact version is bundled unchanged as \texttt{sources/210013/source0.tex}.
\begin{lemma}\label{lemma: cammini}
Let $A\subset \mathbb{R}^{d}$ be a connected set which is finite union of open rectangles $A=\cup_{h=1}^NA_h$ with $N\geq2$ and such that
\begin{equation*}
    A_h\cap A_{h+1}\cap \s\ZZ^d\neq \emptyset \qquad \text{ for every } h\in\{1,\dots,N-1\}.
\end{equation*}
There exist positive constants $\s_{0}$ and $C$ depending on $A_1,\dots,A_N$ such that it holds 
\begin{align}\label{eq: archi}
    \sum_{1\leq j<\ell\leq N}\sum_{\alpha\in (A_j\setminus A_\ell) \cap\s\ZZ^d}\sum_{\beta \in (A_\ell\setminus A_j)\cap\s\ZZ^d}|u(\alpha)-u(\beta)|^p & \leq CN^{p+1} \sum_{h=1}^{N}\sum_{\alpha\in A_h\cap \s \ZZ^d}\sum_{\beta\in A_h\cap \s \ZZ^d}|u(\alpha)-u(\beta)|^p
\end{align}
    for every $\s<\s_0$ and $u:A\cap\s \ZZ^d\rightarrow \mathbb{R}^{m}$.
\end{lemma}

\begin{proof}
We consider $Q_1,\dots, Q_{N-1}$ open cubes with sides parallel to the coordinate axes with the property that $Q_h\subseteq A_h\cap A_{h+1}$ for every $h\in\{1,\dots, N-1\}$,
\begin{equation*}
   |Q_1|=\dots= |Q_{N-1}|\geq \frac{1}{2}\min\{|A_h\cap A_{h+1}| :  h\in\{1,\dots, N-1\}\}, 
\end{equation*}
and 
\begin{equation*}
   C(\s):= \#(Q_1\cap \s\ZZ^d)=\dots = \#(Q_{N-1}\cap \s\ZZ^d)
\end{equation*}
for every $\s$ small enough. We label the elements of each set $Q_h\cap \s\ZZ^d$ as $\{q_h^1,\dots, q_h^{C(\s)}\}$.

We define surjective functions
\begin{equation*}
   I_{h}: A_h\cap\s\mathbb Z^d \to \{1,\dots, C(\s)\}, \qquad h\in\{1,\dots,N\},
\end{equation*}
such that
\begin{equation*}
    \#(I_h^{-1}(\{k\})) \leq \Bigl\lceil\frac{\#(A_{h}\cap \s \mathbb Z^d)}{C(\s)}\Bigr\rceil
\end{equation*}
for every $h\in\{1,\dots,N\}$ and $k\in\{1,\dots, C(\s)\}$, where we let $\lceil x\rceil$ denote the upper integer part of $x$. Upon assuming that $\s$ is sufficiently small, we can then suppose that
\begin{equation}\label{stima cardinalità}
    \#(I_h^{-1}(\{k\})) \leq 8|A_h|/\min\{|A_h\cap A_{h+1}| :  h\in\{1,\dots, N-1\}\} 
\end{equation}
for every $h\in\{1,\dots,N\}$ and $k\in\{1,\dots, C(\s)\}$. 

Let $1\leq j<\ell\leq N$ and consider points $\alpha\in (A_j\setminus A_\ell)\cap\s\ZZ^d$ and $\beta \in (A_\ell\setminus A_j)\cap\s\ZZ^d$. We define the intermediate nodes that shall allow us to connect these points. We set
 \begin{equation*}
     p_{j+k}:=\begin{cases}
         q_{j+k}^{I_\ell(\beta)} & \text{ if } k \text{ is even,}\\[5pt]
         q_{j+k}^{I_j(\alpha)} & \text{ if } k \text{ is odd,}
     \end{cases}
     \qquad k\in\{0,\dots,\ell-j-1\}.
 \end{equation*}
For sake of notation, it is convenient to set 
 \begin{equation*}
n(j,\ell):= \begin{cases}
        \ell & \text{ if } \ell-j \text{ is even,} \\
     \ell+1 & \text{ if } \ell-j \text{ is odd,}
    \end{cases}   
\end{equation*}
and 
\begin{equation*}
    p_{j-1}:=\alpha, \qquad p_{n(j,\ell)}:=\beta.
\end{equation*}
Now we distinguish two cases:
\begin{itemize}
 \item if $\ell-j$ is even, we have that $p_{\ell-1}=q_{\ell-1}^{I_j(\alpha)}$ and we define
 \begin{equation*}
     \mathcal P(\alpha,\beta):=
        \{\alpha=p_{j-1}, p_j, \dots, p_{\ell-1}, p_\ell=p_{n(j,\ell)}=\beta\}, 
 \end{equation*}
\item if $\ell-j$ is odd, we have that $p_{\ell-1}=q_{\ell-1}^{I_\ell(\beta)}$, then we set $p_\ell:=q_{\ell-1}^{I_j(\alpha)}$ and we define
\begin{equation*}
    \mathcal P(\alpha,\beta):=\{\alpha=p_{j-1}, p_j, \dots, p_{\ell}, p_{\ell+1}=p_{n(j,\ell)}=\beta\} .
\end{equation*}   
 \end{itemize}
Finally, we connect $\alpha$ and $\beta$ using the path obtained connecting through segments the points of $\mathcal P(\alpha,\beta)$ in the established order.

If $\ell-j$ is even, the path associated with $\mathcal P(\alpha,\beta)$ is made up of at most $\ell-j+1$ segments and 
\begin{equation*}
 [p_h,p_{h+1}]\subset A_{h+1} \quad \text{ for every } h\in\{j-1,\dots, \ell-1\}.
\end{equation*}

If $\ell-j$ is odd, the path associated with $\mathcal P(\alpha,\beta)$ is made up of at most $\ell-j+2$ segments and 
\begin{equation*}
  [p_h,p_{h+1}]\subset A_{h+1} \quad\text{ for every } h\in\{j-1,\dots, \ell-1\}, \qquad [p_{\ell},p_{\ell+1}]\subset A_\ell.
\end{equation*}

Let us consider a segment $[\lambda,\eta]\subset A_{h}$ for some $h\in\{1,\dots,N\}$. We let $S(\lambda,\eta)$ denote the set of pairs $(\alpha,\beta)\in(A_j\setminus A_\ell)\cap\s\ZZ^d\times (A_\ell\setminus A_j)\cap\s\ZZ^d, 1\leq j<\ell\leq N,$ such that $\lambda$ and $\eta$ are two consecutive points belonging to the path associated with $\mathcal{P}(\alpha,\beta)$. 

To estimate the cardinality of $S(\lambda,\eta)$, let us assume that $\ell-j$ is even and consider $(\alpha,\beta)\in S(\lambda,\eta)$.  Upon exchanging the roles of $\lambda$ and $\eta$, we have the following possibilities:
\begin{itemize}
    \item[(i)] $\lambda=\alpha=p_{j-1}, \eta=p_j$,
     \item[(ii)] $  \lambda=p_{\ell-1}, \eta=\beta=p_\ell$,
      \item[(iii)] $\lambda=p_m,  \eta=p_{m+1}$ for some $m\in\{j,\dots,\ell-2\}$.
\end{itemize}
In the first case $\alpha$ is uniquely determined and $\eta\in Q_j$, so that $\eta=q_j^{k}$ for a certain $k\in\{1,\dots, C(\s)\}$. Hence, we have that 
\begin{equation*}
q_j^k=\eta=p_j=q^{I_\ell(\beta)}_j.    
\end{equation*}
Therefore, $\beta\in I_\ell^{-1}(\{k\})$ and, by \eqref{stima cardinalità}, the number of possible pairs $(\alpha,\beta)$ is less than or equal to 
\begin{equation*}
    \#(I_\ell^{-1}(\{k\})) \leq 8|A_\ell|/\min\{|A_h\cap A_{h+1}| :  h\in\{1,\dots, N-1\}\}.
\end{equation*}
In the second case $\beta$ is uniquely determined and $\lambda\in Q_{\ell-1}$, so that $\lambda=q_{\ell-1}^{k}$ for a certain $k\in\{1,\dots, C(\s)\}$. Hence, we have that 
\begin{equation*}
    q_{\ell-1}^{k}= \lambda=p_{\ell-1}=q^{I_j(\alpha)}_{\ell-1}.
\end{equation*}
Therefore, $\alpha\in I_j^{-1}(\{k\})$ and the number of possible pairs $(\alpha,\beta)$ is less than or equal to 
\begin{equation*}
    \#(I_j^{-1}(\{k\})) \leq  8|A_j|/\min\{|A_h\cap A_{h+1}| :  h\in\{1,\dots, N-1\}\}.
\end{equation*}
In the third case, since $\lambda\in Q_m$ and $\eta\in Q_{m+1}$, we have that $\lambda=q_m^k$ and $\eta=q_{m+1}^{k'}$ for some $k,k'\in\{1,\dots, C(\s)\}$, and then
\begin{equation*}
  \begin{cases}
      q_m^k=q_m^{I_\ell(\beta)}, \\[3pt]
      q_{m+1}^{k'}=q_{m+1}^{I_j(\alpha)},
  \end{cases}\qquad \text{ or } \qquad  
  \begin{cases}
q_m^k=q_{m}^{I_j(\alpha)}, \\[3pt]
q_{m+1}^{k'}=q_{m+1}^{I_\ell(\beta)}.
  \end{cases}
\end{equation*}
Therefore,
\begin{equation*}
\begin{cases}
   \alpha\in I_j^{-1}(\{k'\}), \\
   \beta \in I_\ell^{-1}(\{k\}), 
\end{cases}
    \qquad \text{ or } \qquad
    \begin{cases}
        \alpha\in I_j^{-1}(\{k\}),\\ \beta \in I_\ell^{-1}(\{k'\}),
    \end{cases} 
\end{equation*}
and the number of possible pairs $(\alpha,\beta)$ in this case is less than or equal to 
\begin{align*}
   & \max\{\#(I_j^{-1}(\{k'\}))\times \#(I_\ell^{-1}(\{k\})), \#(I_j^{-1}(\{k\}))\times \#(I_\ell^{-1}(\{k'\}))\} \\ &\leq64|A_j||A_\ell|/\min\{|A_h\cap A_{h+1}| :  h\in\{1,\dots, N-1\}\}^2.
\end{align*}

The case $\ell-j$ odd is analogous, provided that the additional segment $[p_{\ell-1},p_\ell]$ is also considered.
For this segment, upon exchanging the roles of $\lambda$ and $\eta$, one has
\[
\lambda=p_{\ell-1}=q_{\ell-1}^{I_\ell(\beta)},\qquad
\eta=p_\ell=q_{\ell-1}^{I_j(\alpha)}.
\]
Therefore, the number of admissible pairs $(\alpha,\beta)$ is bounded by
\[
\#I_\ell^{-1}(\{k\})\times\# I_j^{-1}(\{k'\})
\leq
64|A_j||A_\ell|/
\min\{|A_h\cap A_{h+1}|:h \in \{1,\dots,N-1\}\}^2,
\]
which is the same estimate obtained for (iii) in the case $\ell-j$ is even. All other segments are treated as in the even case.

Letting $j,\ell$ vary in $\{1,\dots, N\}$, we obtain that
\begin{align} \notag
    \#S(\lambda,\eta) & \leq 2(64+8) \sum_{1\leq j<\ell\leq N} |A_j||A_\ell|/\min\{|A_h\cap A_{h+1}| :  h\in\{1,\dots, N-1\}\}^2 \\ \label{stima archi}
    & \leq N^2 144 \Bigl(\frac{\max\{|A_h|: h\in\{1,\dots, N\}\}}{\min\{|A_h\cap A_{h+1}| :  h\in\{1,\dots, N-1\}\}}\Bigr)^2=:N^2C.
\end{align}

By Jensen's inequality we have 
\begin{align*}
    |u(\alpha)-u(\beta)|^p &= \Bigl|\sum_{h=j-1}^{n(j,\ell)-1}u(p_{h+1})-u(p_h)\Bigr|^p \\
    &\leq (n(j,\ell)-j+1)^{p-1} \sum_{h=j-1}^{n(j,\ell)-1}|u(p_{h+1})-u(p_h)|^p \leq (2N)^{p-1}\sum_{h=j-1}^{n(j,\ell)-1}|u(p_{h+1})-u(p_h)|^p.
\end{align*}
Recalling that $[p_h,p_{h+1}]\subset A_{h+1}$ for every $h\in\{j-1,\dots, \ell-1\}$ and $[p_{n(j,\ell)-1},p_{n(j,\ell)}]\subset A_\ell$, and using \eqref{stima archi}, we sum over $\alpha$ and $\beta$ and $1\leq j<\ell\leq N$ to get 
\begin{multline*}
     \sum_{1\leq j<\ell\leq N}\sum_{\alpha\in (A_j\setminus A_\ell) \cap\s\ZZ^d}\sum_{\beta \in (A_\ell\setminus A_j)\cap\s\ZZ^d}|u(\alpha)-u(\beta)|^p \\ \leq 2^{p-1}N^{p+1}C \sum_{h=1}^{N}\sum_{\lambda\in A_h\cap \s \ZZ^d}\sum_{\eta\in A_h\cap \s \ZZ^d}|u(\lambda)-u(\eta)|^p,
\end{multline*}
which concludes the proof.
\end{proof}

\end{document}
